\section{The MINERvA measurement-dimensionality survey}
\label{app:landscape}

The classification behind Fig.~\ref{fig:landscape} records the surveyed MINERvA
differential cross-section measurements through the listed 2025 entries, the
number of kinematic variables each headline result unfolds simultaneously,
and the reference. Additional recent measurements are discussed below; this
inventory is not an exhaustive census through 2026.

{\small
\begin{longtable}{@{}clcl@{}}
  \caption{MINERvA differential cross-section survey used for the landscape figure: the number of
  kinematic variables unfolded simultaneously (Dim.), classified from each
  analysis's headline differential result. Channels: CC-inclusive (incl.),
  CCQE/QE-like incl.\ transverse kinematic imbalance (QE), pion production CC
  and coherent ($\pi$), deep-inelastic / kaon (DIS), electron-neutrino
  ($\nu_e$), nuclear-target ratio (ratio).}
  \label{tab:landscape}\\
  \toprule
  Year & Channel; observables & Dim. & Ref. \\
  \midrule
  \endfirsthead
  \toprule
  Year & Channel; observables & Dim. & Ref. \\
  \midrule
  \endhead
  \midrule \multicolumn{4}{r}{\emph{continued\ldots}}\\ \endfoot
  \bottomrule
  \endlastfoot
  2013 & QE; $\nu_\mu$ CCQE $\pT,\pPar$            & 2 & \cite{MINERvA:2013qe} \\
  2013 & QE; $\bar\nu_\mu$ CCQE $\pT,\pPar$        & 2 & \cite{MINERvA:2013qebar} \\
  2014 & ratio; CC on C/Fe/Pb to CH               & 1 & \cite{MINERvA:2014ratio} \\
  2014 & $\pi$; coherent $\pi^\pm$ (total \& diff.) & 1 & \cite{MINERvA:2014coh} \\
  2015 & $\pi$; CC $1\pi^+$ ($T_\pi$, $\theta_\pi$) & 1 & \cite{MINERvA:2015ydy} \\
  2015 & QE; $\mu$+p final states                 & 1 & \cite{MINERvA:2015mup} \\
  2015 & $\pi$; $\bar\nu_\mu$ single $\pi^0$       & 1 & \cite{MINERvA:2015pi0} \\
  2016 & $\nu_e$; CCQE-like on CH                  & 2 & \cite{MINERvA:2016nueqe} \\
  2016 & $\pi$; $\nu_\mu/\bar\nu_\mu$ pion prod.   & 2 & \cite{MINERvA:2016pion} \\
  2016 & DIS; partonic nuclear effects            & 2 & \cite{MINERvA:2016dis} \\
  2016 & DIS; CC $K^+$ production                  & 1 & \cite{MINERvA:2016kaon} \\
  2017 & $\pi$; CC single $\pi^0$                  & 2 & \cite{MINERvA:2017pi0} \\
  2017 & QE; nuclear-dependence (multi-target)     & 1 & \cite{MINERvA:2017Adep} \\
  2018 & QE; $\bar\nu_\mu$ QE-like $\pT,\pPar$     & 2 & \cite{MINERvA:2018hba} \\
  2018 & QE; $\mu$-p mesonless TKI                 & 2 & \cite{MINERvA:2018tcb} \\
  2018 & $\pi$; coherent $\pi^\pm$ on C            & 2 & \cite{MINERvA:2018coh} \\
  2019 & QE; quasielastic-like $\nu_\mu$           & 2 & \cite{MINERvA:2019qe} \\
  2019 & $\pi$; $\bar\nu_\mu$ single $\pi^-$       & 2 & \cite{MINERvA:2019pim} \\
  2020 & incl.; CC-incl $\langle E_\nu\rangle\!\sim\!3.5$ GeV $\pT,\pPar$ & 2 & \cite{MINERvA:2020jwd} \\
  2020 & QE; nuclear binding \& TKI               & 2 & \cite{MINERvA:2020cgh} \\
  2020 & QE; high-statistics QE-like              & 2 & \cite{MINERvA:2020qehs} \\
  2020 & $\pi$; CC $\pi^0$ nuclear effects        & 1 & \cite{MINERvA:2020pi0} \\
  2021 & incl.; CC-incl $\langle E_\nu\rangle\!\sim\!6$ GeV $\pT,\pPar$ \textbf{(target)} & 2 & \cite{MINERvA:2021owq} \\
  2022 & QE; QE-like $\pT,\pPar,\Sigma T_p$ (triple-diff.) & 3 & \cite{MINERvA:2022qe} \\
  2022 & incl.; CC-incl low momentum transfer     & 1 & \cite{MINERvA:2022incl} \\
  2023 & $\pi$; CC $1\pi^+$ (multi-target)        & 2 & \cite{MINERvA:2023ccpi} \\
  2023 & $\pi$; coherent $\pi^+$ (multi-target)   & 1 & \cite{MINERvA:2023coh} \\
  2023 & QE; $\bar\nu_\mu$ multi-neutron low-$\Eavail$ & 1 & \cite{MINERvA:2023lowEavail} \\
  2024 & $\nu_e$; low-$Q^2$ ($\Eavail$, $\pT$)    & 2 & \cite{MINERvA:2023avz} \\
  2025 & QE; $A$-dependence TKI (multi-target)    & 2 & \cite{MINERvA:2025tki} \\
  \midrule
  2026 & \textbf{this work}; unbinned $\pT,\pPar$            & 2 & --- \\
  2026 & \textbf{this work}; unbinned $\pT,\pPar,\Eavail$    & \textbf{3} & --- \\
  2026 & \textbf{this work}; unbinned $\pT,\pPar,\Eavail,q_{3}$ & \textbf{4} & --- \\
  2026 & \textbf{this work}; unbinned $\pT,\pPar,\Eavail,q_{3},W$ & \textbf{5} & --- \\
\end{longtable}
}

\paragraph{Additional recent measurements.}
The low-/medium-energy comparison of quasielastic-like cross sections reports
three simultaneous axes, $\pT$, $\pPar$, and available recoil energy
\cite{MINERvA:2026qeleme}; it is a direct dimensional precedent with a different
signal from our inclusive result. The shallow-inelastic measurement reports
one-dimensional neutrino and antineutrino cross sections
\cite{MINERvA:2025sis}. Further measurements report single-differential
charged-pion production extending to zero pion kinetic energy
\cite{MINERvA:2026pion} and inclusive antineutrino cross sections on several
targets as functions of $\pT$~\cite{MINERvA:2026antinu}. The number of measured
observables is not the number of axes unfolded simultaneously. These additions
provide channel-specific context and are not plotted in Fig.~\ref{fig:landscape}.
